\documentclass[sigconf,review,anonymous]{acmart}
\usepackage{graphicx}

% Draft only: author, venue, and publication rights metadata are pending.
\setcopyright{none}
\settopmatter{printacmref=false,printccs=false}
\copyrightyear{}
\acmConference[]{}{}{}
\acmDOI{}
\acmISBN{}
\title{Constructing Peace in UN Discourse on the Palestine Question}
\author{Author information withheld for review}

\begin{document}
\begin{abstract}
\emph{Placeholder: Write the abstract after human validation and interpretation.}
\end{abstract}
\keywords{United Nations, Palestine question, peace discourse, Galtung, content analysis}
\maketitle
% No venue metadata exists yet, so omit ACM's otherwise empty conference header.
\pagestyle{plain}

\section{Introduction}
\emph{Placeholder: Develop the research problem, contribution, and motivation.}

\section{Related Work and Theoretical Framework}
\emph{Placeholder: Add verified scholarship on peace theory, diplomatic discourse,
and computational content analysis.}

\section{Research Questions}
\begin{enumerate}
  \item How is peace constructed in UN member-state discourse on the Palestine question
  through Galtungian dimensions of peace, violence, relational structure, domination,
  and conflict transformation?
  \item How does the relative emphasis on these dimensions vary across member states?
  \item What recurring Galtungian peace profiles emerge among member states, and which
  states exhibit similar or contrasting configurations?
\end{enumerate}

\section{Corpus and Measurement}
\subsection{Meeting and Speaker Selection}
The meeting search covers 7 October 2023 through 31 December 2025, inclusive.
A meeting enters the statistical corpus only when the collected UN meeting manifest
records an agenda item naming Palestine. This meeting-level rule does not depend
on incidental Palestine or Gaza mentions in speeches; it also excludes Gaza-focused
debates held under other agenda items. The latest qualifying meeting in the
collection was on 16 December 2025 \cite{UNSC10068}. A meeting on 29 December
2025 about Israel's recognition of Somaliland is excluded despite incidental
Palestine and Gaza references \cite{UNSC10084}.

The manifest identifies 112 meetings satisfying the date and agenda rules.
The selected annotation runs cover 93 of them; the other 19, mostly General
Assembly records, are listed in the meeting-scope audit. The results therefore
describe the annotated subset, not all eligible UN meetings. Speeches are
segmented into paragraphs. The main population consists of substantive,
nonprocedural paragraphs spoken by member-state representatives. UN officials
and observer cases, including Palestine, are excluded from the main state
comparison; Israel remains a member-state case. RQ1 includes Security Council
and General Assembly records and contains 14,463 paragraphs from 1,751 speeches.
RQ2 and RQ3 use the Security Council subset, with 14,070 paragraphs before
country eligibility thresholds are applied.
% Sources: audit/meeting_scope.csv and reports/provisional_report.md.

\subsection{Paragraph Annotation}
Five theory-guided codebooks measure negative and positive peace, forms of
violence, relationships between parties, structural domination, and conflict
transformation. The analysis selects one completed full annotation run for each:
version 1 for positive/negative peace, violence, imperialism, and transcend,
and version 2 for relations. A paragraph may receive multiple codes, each
with an evidence span. Repeated evidence for the same paragraph, run, dimension,
and code is counted once. Actor attribution in the violence codebook is
supplementary and excluded from the core country profile.

For each code, a successfully completed annotation job makes the paragraph
evaluable. A completed job without that code contributes an observed zero;
an absent, failed, or incomplete job remains missing. Violence attribution is
evaluable only when a violence type was coded. Thus each code's prevalence
denominator comprises its evaluable paragraphs, rather than every corpus
paragraph. The annotation database was read without altering or rerunning
the model judgments.
% Sources: config/research.yaml and audit/audit_report.md.

\subsection{Human Assessment Status}
Two blinded reviewer packs are prepared from scoped, nonprocedural paragraphs.
For each of five codebooks, the design contains 100 random probability-sample
tasks and 50 purposive diagnostic tasks, giving 750 tasks over 730 distinct
paragraphs. The diagnostic tasks emphasize rare, multi-label, and lower-confidence
cases and will be assessed separately from the probability sample. Reviewers
have not returned judgments, so agreement, adjudication, and measurement-error
estimates remain pending. All findings in this draft are provisional and unvalidated.
% Source: validation/coordinator_only/sampling_manifest.json.

\section{Methods}
\subsection{RQ1: Shared Repertoire}
For each code, paragraph prevalence is the fraction of evaluable member-state
paragraphs carrying that code. Ninety-five percent intervals come from 500
bootstrap draws of speeches, carrying their paragraphs together. A complementary
speech-weighted measure averages within-speech code shares across speeches.
Pairwise co-occurrence uses only paragraphs evaluable for both codes; joint
prevalence, phi correlation, and lift describe their overlap. Reproducibly
selected high-confidence and boundary passages are reserved for close reading,
not treated as qualitative validation by themselves.

\subsection{RQ2: Cross-Country Variation}
Country comparisons use Security Council speeches. Primary eligibility requires
at least 100 unique paragraphs, 15 speeches, and 20 evaluated paragraphs for
every core codebook. Sensitivity thresholds are 50 paragraphs, 10 speeches,
and 10 evaluated paragraphs per core codebook. We report paragraph and
speech averages separately.
Each core code is modeled as a binary paragraph outcome when at least 30
positive and 30 negative observations are available overall.

The primary mixed model contains country and calendar-month terms, log
paragraph word count, and random intercepts for speech and meeting. It is
fitted with mean-field variational Bayes. Country probabilities are predicted
over the same observed month and length distribution for every country, with
random intercepts fixed at zero. Approximate 95\% intervals use 400 draws
of the fitted fixed-effect approximation. Omnibus country Wald tests and
Benjamini--Hochberg adjustments are approximate. A planned contrast
compares Israel with the mean for other eligible states. Refits omitting
length or replacing month with quarter check model specification. Nonconverged
refits remain visible in diagnostics and are omitted from sensitivity
predictions.

\subsection{RQ3: Country Profiles}
RQ3 arranges converged adjusted country probabilities from five core theory
blocks in a country-by-code matrix. Actor attribution is excluded and missing
values are not filled with zero. Each code is standardized across countries
and divided by the square root of the number of retained codes in its block,
so larger codebooks do not dominate. Principal component analysis describes
continuous variation. Ward clustering with Euclidean distance is examined
for two through five profiles using silhouette and group-size diagnostics;
neutral profile identifiers are exploratory, not political labels.

Two distinct checks probe the geometry without hundreds of mixed-model refits.
First, 400 draws from the fitted mean-field fixed-effect approximations are
reclustered. This measures uncertainty conditional on the fitted models
without full posterior covariance or speech refitting. Second, speeches are
resampled 400 times within country to recompute and recluster
\emph{descriptive} profiles. This measures sensitivity to speech composition,
not uncertainty in adjusted model predictions. Co-clustering rates, adjusted
Rand indices, and best-match Jaccard scores summarize the checks separately.
Neither check validates stable political groups or replaces a full model-refit
bootstrap.

% Keep each visual-only Results subsection with its own tables and figures.
\clearpage
\section{Results}
\subsection{RQ1: Galtungian Repertoire}
\input{../research_outputs/latex/paper/rq1_prevalence.tex}
\begin{figure*}[tbp]
\centering
\includegraphics[width=0.88\textwidth]{../research_outputs/rq1_provisional/prevalence.png}
\caption{Paragraph prevalence by Galtung code with speech-resampled 95\% intervals.
Provisional and unvalidated.}
\Description{Horizontal bars show estimated prevalence for 21 core codes,
grouped by codebook, with interval marks at each bar end. Violence codes VD
and VS are among the longest bars; relation code AB is near zero.}
\label{fig:rq1-prevalence}
\end{figure*}
\begin{figure*}[tbp]
\centering
\includegraphics[width=0.80\textwidth]{../research_outputs/rq1_provisional/cooccurrence_phi.png}
\caption{Pairwise code association on jointly evaluable paragraphs.
Provisional and unvalidated.}
\Description{Square heatmap with Galtung code labels on both axes. A diverging
color scale encodes pairwise phi correlation across the matrix.}
\label{fig:rq1-cooccurrence}
\end{figure*}
\clearpage

\subsection{RQ2: Cross-Country Variation}
\input{../research_outputs/latex/paper/rq2_models.tex}
\clearpage

\subsection{RQ3: Similarities and Contrasts}
\input{../research_outputs/latex/paper/rq3_clusters.tex}
\input{../research_outputs/latex/paper/rq3_light_stability.tex}
\begin{figure*}[tbp]
\centering
\includegraphics[width=0.95\textwidth]{../research_outputs/rq3_provisional/fig_cluster_diagnostics.png}
\caption{Ward clustering diagnostics for adjusted country profiles.
Provisional and unvalidated.}
\Description{Three panels compare mean silhouette width for two through five
clusters, smallest-cluster size and largest-cluster share, and individual
country silhouette widths for the two-cluster solution.}
\label{fig:rq3-cluster-diagnostics}
\end{figure*}
\begin{figure*}[tbp]
\centering
\includegraphics[width=0.94\textwidth]{../research_outputs/rq3_provisional/dendrogram.png}
\caption{Ward dendrogram of standardized, block-scaled adjusted country profiles.
Provisional and unvalidated.}
\Description{Hierarchical tree with country names at the leaves and branch
height representing Ward linkage distance between adjusted profiles.}
\label{fig:rq3-dendrogram}
\end{figure*}
\begin{figure*}[tbp]
\centering
\includegraphics[width=0.94\textwidth]{../research_outputs/rq3_provisional/fig_knn_similarity_network.png}
\caption{Supplementary local-similarity network of adjusted country profiles.
Provisional and unvalidated.}
\Description{Countries appear as labeled nodes positioned by the first two
principal components. Edges connect nearest neighbors based on distances in
the complete standardized, block-scaled profile.}
\label{fig:rq3-network}
\end{figure*}
\clearpage

\section{Discussion}
\emph{Placeholder: Integrate the three research questions after human assessment
and close reading of selected passages.}

\section{Limitations and Research Ethics}
\emph{Placeholder: Develop the implications of agenda selection, incomplete
annotations, speech opportunity, model approximation, translation, and
measurement error.}

\section{Conclusion}
\emph{Placeholder: Write conclusions after validation and interpretation.}

\bibliographystyle{ACM-Reference-Format}
\bibliography{references}
\end{document}
